\scanpage{72}
\mspage{10}
\uline{NB 1} Cette fonction est unique, dû au fait
que \deleted{\unsure{les|le}} le groupe des automorphismes \deleted{conformes}
holomorphes de $\mathbb D$ agit de façon simplement
transitive sur l'ens des
triples \deleted{\unsure{d'él.|ord}} de pts distincts $(r,s,t)$ \inserted{sur $\mathbb{S}^{1}=\partial\mathbb D$} qui se
suivent pour l'orientation canonique de
$\deleted{\partial\mathbb D}\mathbb S^1$. \deleted{\unsure{D'où|Donc} \unsure{il}}

\uline{NB 2} La condition 3) implique qu'on a un
\unsure{homéo|bijection} canonique
\[
\varphi:\beta=(\widehat{1,\zeta})\simeq\bar\beta=(\widehat{1,\bar\zeta})
\]
tel que $\varphi(1)=1$, $\varphi(\zeta)=\bar\zeta$).

B) Cette application $\varphi$ est rectifiable. Inversement,
\deleted{toute} tout \deleted{application} \inserted{homéo} rectifiable $\varphi:\beta\xrightarrow{\sim}\bar\beta$
provient d'une situation comme en A.
De façon plus précise, \deleted{on peut \unsure{supp.} \inserted{\deleted{il existe}}}
(quitte à remplacer $\bar V$ par $\mathbb D$ par le th. d'uniformisation A)
simplifié sans segment intérieur $I$,
et normaliser la situation en \deleted{\unsure{plaçant|plaçan}} ramenant
$a$ en $0$, $b$ en $1$) \deleted{il existe \unsure{un|une}} Il existe
une application continue
\[
f=f_\varphi:\mathbb D\mathrel{\unsure{-|\to}}\mathbb D
\]
telle que

1) $f\mid\mathbb D^\circ$ soit \deleted{holo} injection holomorphe
% archive sheet number: 74
